\section{Embedding $\mathbb{Z}$ into $\mathbb{Q}$}
\label{sec:embedding-z-into-q}

\begin{topicbox}{Embedding $\mathbb{Z}$ into $\mathbb{Q}$}
\begin{exposition}
This section fixes the canonical embedding of $\mathbb{Z}$ into $\mathbb{Q}$ and records that it is injective, structure-preserving, and order-reflecting. The map is the bridge that lets elements of $\mathbb{Z}$ be regarded as elements of $\mathbb{Q}$ after identification.
\end{exposition}
\end{topicbox}
\begin{definitionbox}{Definition (Embedding $\mathbb{Z}$ into $\mathbb{Q}$)}
\begin{definition}[Embedding $\mathbb{Z}$ into $\mathbb{Q}$]
\label{def:embedding-z-into-q}
The \emph{canonical embedding} $\iota:\mathbb{Z}\to \mathbb{Q}$ is defined by
\[
\iota(z)=\left[(z,1)\right] \quad(\text{the rational represented by } z/1).
\]
\end{definition}
\end{definitionbox}
\begin{remark*}[Standard quantified statement]
\[
\iota:\mathbb{Z}\to \mathbb{Q},\qquad \iota(z)=\left[(z,1)\right] \quad(\text{the rational represented by } z/1).
\]
\end{remark*}
\begin{remark*}[Interpretation]
The map sends each element of $\mathbb{Z}$ to its natural representative in $\mathbb{Q}$. The next result certifies that this representative choice preserves all arithmetic and order, so $\mathbb{Z}$ may be identified with its image in $\mathbb{Q}$.
\end{remark*}
\begin{dependencies}
\begin{itemize}
  \item \hyperref[def:rational-numbers]{Construction of $\mathbb{Q}$ as fraction classes of integers}
  \item \hyperref[def:number-system-embedding]{Embedding of Number Systems}
\end{itemize}
\end{dependencies}
\begin{theorembox}{Theorem ($\mathbb{Z}$ Embeds in $\mathbb{Q}$)}
\begin{theorem}[$\mathbb{Z}$ Embeds in $\mathbb{Q}$]
\label{thm:z-embeds-in-q}
The canonical map $\iota:\mathbb{Z}\to \mathbb{Q}$ is an embedding of ordered arithmetic systems: it is injective, preserves $0$, $1$, addition, and multiplication, and reflects the order.
\hyperref[prf:z-embeds-in-q]{Go to proof.}
\end{theorem}
\end{theorembox}
\begin{remark*}[Standard quantified statement]
\[
\iota:\mathbb{Z}\to \mathbb{Q} \text{ is injective, structure-preserving, and order-reflecting.}
\]
\end{remark*}
\begin{remark*}[Interpretation]
Once this theorem is proved, the identification theorem of \S\ref{sec:embeddings-as-structure-preserving-maps} applies, and $\mathbb{Z}$ may be treated as the substructure $\iota(\mathbb{Z})\subseteq \mathbb{Q}$.
\end{remark*}
\begin{dependencies}
\begin{itemize}
  \item \hyperref[def:rational-numbers]{Construction of $\mathbb{Q}$ as fraction classes of integers}
  \item \hyperref[def:embedding-z-into-q]{Embedding $\mathbb{Z}$ into $\mathbb{Q}$}
  \item \hyperref[thm:embedding-identifies-with-image]{Embeddings Identify a System with Its Image}
\end{itemize}
\end{dependencies}
